\hypertarget{fichiers}{%
\section{Fichiers}\label{fichiers}}

\hypertarget{persistance}{%
\subsection{Persistance}\label{persistance}}

La plupart des programmes que nous avons vus jusqu\textquotesingle à
présent sont éphémères dans le sens où ils s\textquotesingle exécutent
pendant un court moment et produisent un certain résultat, mais
lorsqu\textquotesingle ils se terminent, leurs données disparaissent. Si
vous relancez le programme, il repart sur une feuille blanche{[}cite:
1{]}.

D\textquotesingle autres programmes sont persistants : ils
s\textquotesingle exécutent longtemps (ou en continu) ; ils conservent
au moins une partie de leurs données dans un stockage permanent (un
disque dur, par exemple) ; et s\textquotesingle ils
s\textquotesingle arrêtent et redémarrent, ils reprennent là où ils
s\textquotesingle étaient arrêtés. Des exemples de programmes
persistants sont les systèmes d\textquotesingle exploitation, qui
s\textquotesingle exécutent à peu près chaque fois qu\textquotesingle un
ordinateur est allumé, et les serveurs web, qui fonctionnent en
permanence en attendant que des requêtes arrivent sur le réseau{[}cite:
1{]}.

L\textquotesingle un des moyens les plus simples pour les programmes de
conserver leurs données consiste à lire et à écrire des fichiers texte.
Nous avons déjà vu des programmes qui lisent des fichiers texte ; dans
ce chapitre, nous verrons des programmes qui en écrivent{[}cite: 1{]}.

Une alternative est de stocker l\textquotesingle état du programme dans
une base de données. Dans ce chapitre, je présenterai une base de
données simple et un module, \texttt{pickle}, qui facilite le stockage
des données de programmes{[}cite: 1{]}.

\hypertarget{lecture-et-uxe9criture}{%
\subsection{Lecture et écriture}\label{lecture-et-uxe9criture}}

Un fichier texte est une séquence de caractères stockée sur un support
permanent comme un disque dur, une mémoire flash ou un cédérom. Nous
avons vu comment ouvrir et lire un fichier à la section 9.1{[}cite:
1{]}.

Pour écrire dans un fichier, vous devez l\textquotesingle ouvrir avec le
mode \texttt{\textquotesingle{}w\textquotesingle{}} comme second
paramètre{[}cite: 1{]} :

\begin{Shaded}
\begin{Highlighting}[]
\OperatorTok{\textgreater{}\textgreater{}\textgreater{}}\NormalTok{ fout }\OperatorTok{=} \BuiltInTok{open}\NormalTok{(}\StringTok{\textquotesingle{}output.txt\textquotesingle{}}\NormalTok{, }\StringTok{\textquotesingle{}w\textquotesingle{}}\NormalTok{)}
\OperatorTok{\textgreater{}\textgreater{}\textgreater{}} \BuiltInTok{print}\NormalTok{ fout}
\OperatorTok{\textless{}}\BuiltInTok{open} \BuiltInTok{file} \StringTok{\textquotesingle{}output.txt\textquotesingle{}}\NormalTok{, mode }\StringTok{\textquotesingle{}w\textquotesingle{}}\NormalTok{ at }\BaseNTok{0xb7eb2410}\OperatorTok{\textgreater{}}
\end{Highlighting}
\end{Shaded}

Si le fichier existe déjà, l\textquotesingle ouvrir en mode écriture
efface les anciennes données et repart à neuf, alors soyez prudent ! Si
le fichier n\textquotesingle existe pas, un nouveau est créé{[}cite:
1{]}.

La méthode \texttt{write} place des données dans le fichier{[}cite: 1{]}
:

\begin{Shaded}
\begin{Highlighting}[]
\OperatorTok{\textgreater{}\textgreater{}\textgreater{}}\NormalTok{ line1 }\OperatorTok{=} \StringTok{"This here\textquotesingle{}s the wattle,}\CharTok{\textbackslash{}n}\StringTok{"}
\OperatorTok{\textgreater{}\textgreater{}\textgreater{}}\NormalTok{ fout.write(line1)}
\end{Highlighting}
\end{Shaded}

Encore une fois, l\textquotesingle objet fichier garde une trace de sa
position, donc si vous appelez \texttt{write} à nouveau, il ajoute les
nouvelles données à la fin{[}cite: 1{]}.

\begin{Shaded}
\begin{Highlighting}[]
\OperatorTok{\textgreater{}\textgreater{}\textgreater{}}\NormalTok{ line2 }\OperatorTok{=} \StringTok{"the emblem of our land.}\CharTok{\textbackslash{}n}\StringTok{"}
\OperatorTok{\textgreater{}\textgreater{}\textgreater{}}\NormalTok{ fout.write(line2)}
\end{Highlighting}
\end{Shaded}

Lorsque vous avez fini d\textquotesingle écrire, vous devez fermer le
fichier{[}cite: 1{]} :

\begin{Shaded}
\begin{Highlighting}[]
\OperatorTok{\textgreater{}\textgreater{}\textgreater{}}\NormalTok{ fout.close()}
\end{Highlighting}
\end{Shaded}

\hypertarget{opuxe9rateur-de-formatage}{%
\subsection{Opérateur de formatage}\label{opuxe9rateur-de-formatage}}

L\textquotesingle argument de \texttt{write} doit être une chaîne de
caractères, donc si nous voulons placer d\textquotesingle autres valeurs
dans un fichier, nous devons les convertir en chaînes. Le moyen le plus
simple d\textquotesingle y parvenir est d\textquotesingle utiliser
\texttt{str}{[}cite: 1{]} :

\begin{Shaded}
\begin{Highlighting}[]
\OperatorTok{\textgreater{}\textgreater{}\textgreater{}}\NormalTok{ x }\OperatorTok{=} \DecValTok{52}
\OperatorTok{\textgreater{}\textgreater{}\textgreater{}}\NormalTok{ fout.write(}\BuiltInTok{str}\NormalTok{(x))}
\end{Highlighting}
\end{Shaded}

Une alternative consiste à utiliser l\textquotesingle opérateur de
formatage, \texttt{\%}. Lorsqu\textquotesingle il est appliqué à des
entiers, \texttt{\%} est l\textquotesingle opérateur modulo. Mais
lorsque le premier opérande est une chaîne, \texttt{\%} est
l\textquotesingle opérateur de formatage{[}cite: 1{]}.

Le premier opérande est la chaîne de format, qui contient une ou
plusieurs séquences de format spécifiant comment le second opérande est
formaté. Le résultat est une chaîne de caractères{[}cite: 1{]}.

Par exemple, la séquence de format
\texttt{\textquotesingle{}\%d\textquotesingle{}} signifie que le second
opérande doit être formaté en tant qu\textquotesingle entier (le « d »
signifie « décimal »){[}cite: 1{]} :

\begin{Shaded}
\begin{Highlighting}[]
\OperatorTok{\textgreater{}\textgreater{}\textgreater{}}\NormalTok{ camels }\OperatorTok{=} \DecValTok{42}
\OperatorTok{\textgreater{}\textgreater{}\textgreater{}} \StringTok{\textquotesingle{}}\SpecialCharTok{\%d}\StringTok{\textquotesingle{}} \OperatorTok{\%}\NormalTok{ camels}
\CommentTok{\textquotesingle{}42\textquotesingle{}}
\end{Highlighting}
\end{Shaded}

Le résultat est la chaîne
\texttt{\textquotesingle{}42\textquotesingle{}}, à ne pas confondre avec
la valeur entière 42{[}cite: 1{]}.

Une séquence de format peut apparaître n\textquotesingle importe où dans
la chaîne, ce qui vous permet d\textquotesingle intégrer une valeur dans
une phrase{[}cite: 1{]} :

\begin{Shaded}
\begin{Highlighting}[]
\OperatorTok{\textgreater{}\textgreater{}\textgreater{}}\NormalTok{ camels }\OperatorTok{=} \DecValTok{42}
\OperatorTok{\textgreater{}\textgreater{}\textgreater{}} \StringTok{\textquotesingle{}I have spotted }\SpecialCharTok{\%d}\StringTok{ camels.\textquotesingle{}} \OperatorTok{\%}\NormalTok{ camels}
\CommentTok{\textquotesingle{}I have spotted 42 camels.\textquotesingle{}}
\end{Highlighting}
\end{Shaded}

S\textquotesingle il y a plus d\textquotesingle une séquence de format
dans la chaîne, le second argument doit être un tuple. Chaque séquence
de format correspond à un élément du tuple, dans
l\textquotesingle ordre{[}cite: 1{]}.

L\textquotesingle exemple suivant utilise
\texttt{\textquotesingle{}\%d\textquotesingle{}} pour formater un
entier, \texttt{\textquotesingle{}\%g\textquotesingle{}} pour formater
un nombre à virgule flottante (ne demandez pas pourquoi), et
\texttt{\textquotesingle{}\%s\textquotesingle{}} pour formater une
chaîne{[}cite: 1{]} :

\begin{Shaded}
\begin{Highlighting}[]
\OperatorTok{\textgreater{}\textgreater{}\textgreater{}} \StringTok{\textquotesingle{}In }\SpecialCharTok{\%d}\StringTok{ years I have spotted }\SpecialCharTok{\%g}\StringTok{ }\SpecialCharTok{\%s}\StringTok{.\textquotesingle{}} \OperatorTok{\%}\NormalTok{ (}\DecValTok{3}\NormalTok{, }\FloatTok{0.1}\NormalTok{, }\StringTok{\textquotesingle{}camels\textquotesingle{}}\NormalTok{)}
\CommentTok{\textquotesingle{}In 3 years I have spotted 0.1 camels.\textquotesingle{}}
\end{Highlighting}
\end{Shaded}

Le nombre d\textquotesingle éléments dans le tuple doit correspondre au
nombre de séquences de format dans la chaîne. De plus, les types des
éléments doivent correspondre aux séquences de format{[}cite: 1{]} :

\begin{Shaded}
\begin{Highlighting}[]
\OperatorTok{\textgreater{}\textgreater{}\textgreater{}} \StringTok{\textquotesingle{}}\SpecialCharTok{\%d}\StringTok{ }\SpecialCharTok{\%d}\StringTok{ }\SpecialCharTok{\%d}\StringTok{\textquotesingle{}} \OperatorTok{\%}\NormalTok{ (}\DecValTok{1}\NormalTok{, }\DecValTok{2}\NormalTok{)}
\PreprocessorTok{TypeError}\NormalTok{: }\KeywordTok{not}\NormalTok{ enough arguments }\ControlFlowTok{for} \BuiltInTok{format}\NormalTok{ string}
\OperatorTok{\textgreater{}\textgreater{}\textgreater{}} \StringTok{\textquotesingle{}}\SpecialCharTok{\%d}\StringTok{\textquotesingle{}} \OperatorTok{\%} \StringTok{\textquotesingle{}dollars\textquotesingle{}}
\PreprocessorTok{TypeError}\NormalTok{: illegal argument }\BuiltInTok{type} \ControlFlowTok{for}\NormalTok{ built}\OperatorTok{{-}}\KeywordTok{in}\NormalTok{ operation}
\end{Highlighting}
\end{Shaded}

Dans le premier exemple, il n\textquotesingle y a pas assez
d\textquotesingle éléments ; dans le second, l\textquotesingle élément
est du mauvais type{[}cite: 1{]}.

L\textquotesingle opérateur de formatage est puissant, mais il peut
s\textquotesingle avérer difficile à utiliser. Vous pouvez en lire
davantage à son sujet sur
\url{http://docs.python.org/2/library/stdtypes.html\#string-formatting\%5Bcite}:
1{]}.

Noms de fichiers et chemins
-\/-\/-\/-\/-\/-\/-\/-\/-\/-\/-\/-\/-\/-\/-\/-\/-\/-\/-\/-\/-\/-\/-\/-\/-\/-

Les fichiers sont organisés en répertoires (également appelés « dossiers
»). Chaque programme en cours d\textquotesingle exécution possède un «
répertoire courant », qui est le répertoire par défaut pour la plupart
des opérations{[}cite: 1{]}.

Par exemple, lorsque vous ouvrez un fichier en lecture, Python le
recherche dans le répertoire courant{[}cite: 1{]}.

Le module \texttt{os} fournit des fonctions pour travailler avec les
fichiers et les répertoires (« os » signifie « operating system »).
\texttt{os.getcwd} renvoie le nom du répertoire courant{[}cite: 1{]} :

\begin{Shaded}
\begin{Highlighting}[]
\OperatorTok{\textgreater{}\textgreater{}\textgreater{}} \ImportTok{import}\NormalTok{ os}
\OperatorTok{\textgreater{}\textgreater{}\textgreater{}}\NormalTok{ cwd }\OperatorTok{=}\NormalTok{ os.getcwd()}
\OperatorTok{\textgreater{}\textgreater{}\textgreater{}} \BuiltInTok{print}\NormalTok{ cwd}
\OperatorTok{/}\NormalTok{home}\OperatorTok{/}\NormalTok{dinsdale}
\end{Highlighting}
\end{Shaded}

\texttt{cwd} signifie « current working directory » (répertoire de
travail courant). Le résultat dans cet exemple est
\texttt{/home/dinsdale}, qui est le répertoire personnel
d\textquotesingle un utilisateur nommé \texttt{dinsdale}{[}cite: 1{]}.

Une chaîne comme \texttt{cwd} qui identifie un fichier est appelée un
\textbf{chemin} (\emph{path}). Un chemin relatif part du répertoire
courant ; un chemin absolu part du répertoire le plus haut du système de
fichiers{[}cite: 1{]}.

Les chemins que nous avons vus jusqu\textquotesingle à présent sont de
simples noms de fichiers, ils sont donc relatifs au répertoire courant.
Pour trouver le chemin absolu d\textquotesingle un fichier, vous pouvez
utiliser \texttt{os.path.abspath}{[}cite: 1{]} :

\begin{Shaded}
\begin{Highlighting}[]
\OperatorTok{\textgreater{}\textgreater{}\textgreater{}}\NormalTok{ os.path.abspath(}\StringTok{\textquotesingle{}memo.txt\textquotesingle{}}\NormalTok{)}
\CommentTok{\textquotesingle{}/home/dinsdale/memo.txt\textquotesingle{}}
\end{Highlighting}
\end{Shaded}

\texttt{os.path.exists} vérifie si un fichier ou un répertoire
existe{[}cite: 1{]} :

\begin{Shaded}
\begin{Highlighting}[]
\OperatorTok{\textgreater{}\textgreater{}\textgreater{}}\NormalTok{ os.path.exists(}\StringTok{\textquotesingle{}memo.txt\textquotesingle{}}\NormalTok{)}
\VariableTok{True}
\end{Highlighting}
\end{Shaded}

S\textquotesingle il existe, \texttt{os.path.isdir} vérifie
s\textquotesingle il s\textquotesingle agit d\textquotesingle un
répertoire{[}cite: 1{]} :

\begin{Shaded}
\begin{Highlighting}[]
\OperatorTok{\textgreater{}\textgreater{}\textgreater{}}\NormalTok{ os.path.isdir(}\StringTok{\textquotesingle{}memo.txt\textquotesingle{}}\NormalTok{)}
\VariableTok{False}
\OperatorTok{\textgreater{}\textgreater{}\textgreater{}}\NormalTok{ os.path.isdir(}\StringTok{\textquotesingle{}music\textquotesingle{}}\NormalTok{)}
\VariableTok{True}
\end{Highlighting}
\end{Shaded}

De même, \texttt{os.path.isfile} vérifie s\textquotesingle il
s\textquotesingle agit d\textquotesingle un fichier{[}cite: 1{]}.

\texttt{os.listdir} renvoie une liste des fichiers (et autres
répertoires) dans le répertoire donné{[}cite: 1{]} :

\begin{Shaded}
\begin{Highlighting}[]
\OperatorTok{\textgreater{}\textgreater{}\textgreater{}}\NormalTok{ os.listdir(cwd)}
\NormalTok{[}\StringTok{\textquotesingle{}music\textquotesingle{}}\NormalTok{, }\StringTok{\textquotesingle{}photos\textquotesingle{}}\NormalTok{, }\StringTok{\textquotesingle{}memo.txt\textquotesingle{}}\NormalTok{]}
\end{Highlighting}
\end{Shaded}

Pour illustrer ces fonctions, l\textquotesingle exemple suivant «
parcourt » (\emph{walks}) un répertoire, affiche les noms de tous les
fichiers et s\textquotesingle appelle lui-même de manière récursive sur
tous les répertoires{[}cite: 1{]}.

\begin{Shaded}
\begin{Highlighting}[]
\KeywordTok{def}\NormalTok{ walk(dirname):}
    \ControlFlowTok{for}\NormalTok{ name }\KeywordTok{in}\NormalTok{ os.listdir(dirname):}
\NormalTok{        path }\OperatorTok{=}\NormalTok{ os.path.join(dirname, name)}

        \ControlFlowTok{if}\NormalTok{ os.path.isfile(path):}
            \BuiltInTok{print}\NormalTok{ path}
        \ControlFlowTok{else}\NormalTok{:}
\NormalTok{            walk(path)}
\end{Highlighting}
\end{Shaded}

\texttt{os.path.join} prend un répertoire et un nom de fichier et les
combine pour former un chemin complet{[}cite: 1{]}.

\textbf{Exercice}

Le module \texttt{os} fournit une fonction appelée \texttt{walk} qui est
similaire à celle-ci mais plus polyvalente. Lisez la documentation et
utilisez-la pour afficher les noms des fichiers dans un répertoire donné
et ses sous-répertoires{[}cite: 1{]}. Solution :
\url{http://thinkpython.com/code/walk.py\%5Bcite}: 1{]}.

\hypertarget{capture-dexceptions}{%
\subsection{Capture
d\textquotesingle exceptions}\label{capture-dexceptions}}

Beaucoup de choses peuvent mal tourner lorsque vous essayez de lire et
d\textquotesingle écrire des fichiers. Si vous essayez
d\textquotesingle ouvrir un fichier qui n\textquotesingle existe pas,
vous obtenez une \texttt{IOError}{[}cite: 1{]} :

\begin{Shaded}
\begin{Highlighting}[]
\OperatorTok{\textgreater{}\textgreater{}\textgreater{}}\NormalTok{ fin }\OperatorTok{=} \BuiltInTok{open}\NormalTok{(}\StringTok{\textquotesingle{}bad\_file\textquotesingle{}}\NormalTok{)}
\PreprocessorTok{IOError}\NormalTok{: [Errno }\DecValTok{2}\NormalTok{] No such }\BuiltInTok{file} \KeywordTok{or}\NormalTok{ directory: }\StringTok{\textquotesingle{}bad\_file\textquotesingle{}}
\end{Highlighting}
\end{Shaded}

Si vous n\textquotesingle avez pas la permission
d\textquotesingle accéder à un fichier{[}cite: 1{]} :

\begin{Shaded}
\begin{Highlighting}[]
\OperatorTok{\textgreater{}\textgreater{}\textgreater{}}\NormalTok{ fout }\OperatorTok{=} \BuiltInTok{open}\NormalTok{(}\StringTok{\textquotesingle{}/etc/passwd\textquotesingle{}}\NormalTok{, }\StringTok{\textquotesingle{}w\textquotesingle{}}\NormalTok{)}
\PreprocessorTok{IOError}\NormalTok{: [Errno }\DecValTok{13}\NormalTok{] Permission denied: }\StringTok{\textquotesingle{}/etc/passwd\textquotesingle{}}
\end{Highlighting}
\end{Shaded}

Et si vous essayez d\textquotesingle ouvrir un répertoire en lecture,
vous obtenez{[}cite: 1{]} :

\begin{Shaded}
\begin{Highlighting}[]
\OperatorTok{\textgreater{}\textgreater{}\textgreater{}}\NormalTok{ fin }\OperatorTok{=} \BuiltInTok{open}\NormalTok{(}\StringTok{\textquotesingle{}/home\textquotesingle{}}\NormalTok{)}
\PreprocessorTok{IOError}\NormalTok{: [Errno }\DecValTok{21}\NormalTok{] Is a directory}
\end{Highlighting}
\end{Shaded}

Pour éviter ces erreurs, vous pourriez utiliser des fonctions comme
\texttt{os.path.exists} et \texttt{os.path.isfile}, mais cela prendrait
beaucoup de temps et de code pour vérifier toutes les possibilités (si «
Errno 21 » est un indice, il y a au moins 21 choses qui peuvent mal
tourner){[}cite: 1{]}.

Il vaut mieux foncer et essayer --- et gérer les problèmes
s\textquotesingle ils surviennent ---, ce qui est précisément ce que
fait l\textquotesingle instruction \texttt{try}. La syntaxe est
similaire à une instruction \texttt{if}{[}cite: 1{]} :

\begin{Shaded}
\begin{Highlighting}[]
\ControlFlowTok{try}\NormalTok{:}
\NormalTok{    fin }\OperatorTok{=} \BuiltInTok{open}\NormalTok{(}\StringTok{\textquotesingle{}bad\_file\textquotesingle{}}\NormalTok{)}
    \ControlFlowTok{for}\NormalTok{ line }\KeywordTok{in}\NormalTok{ fin:}
        \BuiltInTok{print}\NormalTok{ line}
\NormalTok{    fin.close()}
\ControlFlowTok{except}\NormalTok{:}
    \BuiltInTok{print} \StringTok{\textquotesingle{}Something went wrong.\textquotesingle{}}
\end{Highlighting}
\end{Shaded}

Python commence par exécuter la clause \texttt{try}. Si tout se passe
bien, il ignore la clause \texttt{except} et continue. Si une exception
se produit, il sort de la clause \texttt{try} et exécute la clause
\texttt{except}{[}cite: 1{]}.

Gérer une exception avec une instruction \texttt{try}
s\textquotesingle appelle \textbf{capturer} (\emph{catch}) une
exception. Dans cet exemple, la clause \texttt{except} affiche un
message d\textquotesingle erreur qui n\textquotesingle est pas très
utile. En général, capturer une exception vous donne
l\textquotesingle opportunité de résoudre le problème, de réessayer, ou
du moins de quitter le programme proprement{[}cite: 1{]}.

\textbf{Exercice}

Écrivez une fonction appelée \texttt{sed} qui prend comme arguments une
chaîne modèle, une chaîne de remplacement et deux noms de fichiers ;
elle doit lire le premier fichier et écrire le contenu dans le second
fichier (en le créant si nécessaire). Si la chaîne modèle apparaît
n\textquotesingle importe où dans le fichier, elle doit être remplacée
par la chaîne de remplacement{[}cite: 1{]}. Si une erreur se produit
lors de l\textquotesingle ouverture, de la lecture, de
l\textquotesingle écriture ou de la fermeture de fichiers, votre
programme doit capturer l\textquotesingle exception, afficher un message
d\textquotesingle erreur et se terminer{[}cite: 1{]}. Solution :
\url{http://thinkpython.com/code/sed.py\%5Bcite}: 1{]}.

\hypertarget{bases-de-donnuxe9es}{%
\subsection{Bases de données}\label{bases-de-donnuxe9es}}

Une base de données est un fichier organisé pour stocker des données. La
plupart des bases de données sont organisées comme un dictionnaire en ce
sens qu\textquotesingle elles associent des clés à des valeurs. La plus
grande différence est que la base de données se trouve sur le disque (ou
un autre stockage permanent), de sorte qu\textquotesingle elle persiste
après la fin du programme{[}cite: 1{]}.

Le module \texttt{anydbm} fournit une interface pour créer et mettre à
jour des fichiers de bases de données. À titre
d\textquotesingle exemple, je vais créer une base de données qui
contient des légendes pour des fichiers d\textquotesingle images{[}cite:
1{]}.

L\textquotesingle ouverture d\textquotesingle une base de données est
similaire à l\textquotesingle ouverture d\textquotesingle autres
fichiers{[}cite: 1{]} :

\begin{Shaded}
\begin{Highlighting}[]
\OperatorTok{\textgreater{}\textgreater{}\textgreater{}} \ImportTok{import}\NormalTok{ anydbm}
\OperatorTok{\textgreater{}\textgreater{}\textgreater{}}\NormalTok{ db }\OperatorTok{=}\NormalTok{ anydbm.}\BuiltInTok{open}\NormalTok{(}\StringTok{\textquotesingle{}captions.db\textquotesingle{}}\NormalTok{, }\StringTok{\textquotesingle{}c\textquotesingle{}}\NormalTok{)}
\end{Highlighting}
\end{Shaded}

Le mode \texttt{\textquotesingle{}c\textquotesingle{}} signifie que la
base de données doit être créée si elle n\textquotesingle existe pas
déjà. Le résultat est un objet base de données qui peut être utilisé
(pour la plupart des opérations) comme un dictionnaire{[}cite: 1{]}.

Si vous créez un nouvel élément, \texttt{anydbm} met à jour le fichier
de la base de données{[}cite: 1{]} :

\begin{Shaded}
\begin{Highlighting}[]
\OperatorTok{\textgreater{}\textgreater{}\textgreater{}}\NormalTok{ db[}\StringTok{\textquotesingle{}cleese.png\textquotesingle{}}\NormalTok{] }\OperatorTok{=} \StringTok{\textquotesingle{}Photo of John Cleese.\textquotesingle{}}
\end{Highlighting}
\end{Shaded}

Lorsque vous accédez à l\textquotesingle un des éléments,
\texttt{anydbm} lit le fichier{[}cite: 1{]} :

\begin{Shaded}
\begin{Highlighting}[]
\OperatorTok{\textgreater{}\textgreater{}\textgreater{}} \BuiltInTok{print}\NormalTok{ db[}\StringTok{\textquotesingle{}cleese.png\textquotesingle{}}\NormalTok{]}
\NormalTok{Photo of John Cleese.}
\end{Highlighting}
\end{Shaded}

Si vous effectuez une autre affectation sur une clé existante,
\texttt{anydbm} remplace l\textquotesingle ancienne valeur{[}cite: 1{]}
:

\begin{Shaded}
\begin{Highlighting}[]
\OperatorTok{\textgreater{}\textgreater{}\textgreater{}}\NormalTok{ db[}\StringTok{\textquotesingle{}cleese.png\textquotesingle{}}\NormalTok{] }\OperatorTok{=} \StringTok{\textquotesingle{}Photo of John Cleese doing a silly walk.\textquotesingle{}}
\OperatorTok{\textgreater{}\textgreater{}\textgreater{}} \BuiltInTok{print}\NormalTok{ db[}\StringTok{\textquotesingle{}cleese.png\textquotesingle{}}\NormalTok{]}
\NormalTok{Photo of John Cleese doing a silly walk.}
\end{Highlighting}
\end{Shaded}

De nombreuses méthodes de dictionnaires, comme \texttt{keys} et
\texttt{items}, fonctionnent également avec les objets de bases de
données. L\textquotesingle itération avec une instruction \texttt{for}
fonctionne de même{[}cite: 1{]}.

\begin{Shaded}
\begin{Highlighting}[]
\ControlFlowTok{for}\NormalTok{ key }\KeywordTok{in}\NormalTok{ db:}
    \BuiltInTok{print}\NormalTok{ key}
\end{Highlighting}
\end{Shaded}

Comme avec les autres fichiers, vous devez fermer la base de données
lorsque vous avez terminé{[}cite: 1{]} :

\begin{Shaded}
\begin{Highlighting}[]
\OperatorTok{\textgreater{}\textgreater{}\textgreater{}}\NormalTok{ db.close()}
\end{Highlighting}
\end{Shaded}

\hypertarget{suxe9rialisation-pickling}{%
\subsection{Sérialisation (Pickling)}\label{suxe9rialisation-pickling}}

Une limitation de \texttt{anydbm} est que les clés et les valeurs
doivent être des chaînes de caractères. Si vous essayez
d\textquotesingle utiliser un autre type, vous obtenez une
erreur{[}cite: 1{]}.

Le module \texttt{pickle} peut aider. Il traduit presque
n\textquotesingle importe quel type d\textquotesingle objet en une
chaîne adaptée au stockage dans une base de données, puis retransforme
les chaînes en objets{[}cite: 1{]}.

\texttt{pickle.dumps} prend un objet en paramètre et renvoie une
représentation sous forme de chaîne (« dumps » est
l\textquotesingle abréviation de « dump string »){[}cite: 1{]} :

\begin{Shaded}
\begin{Highlighting}[]
\OperatorTok{\textgreater{}\textgreater{}\textgreater{}} \ImportTok{import}\NormalTok{ pickle}
\OperatorTok{\textgreater{}\textgreater{}\textgreater{}}\NormalTok{ t }\OperatorTok{=}\NormalTok{ [}\DecValTok{1}\NormalTok{, }\DecValTok{2}\NormalTok{, }\DecValTok{3}\NormalTok{]}
\OperatorTok{\textgreater{}\textgreater{}\textgreater{}}\NormalTok{ pickle.dumps(t)}
\CommentTok{\textquotesingle{}(lp0}\CharTok{\textbackslash{}n}\CommentTok{I1}\CharTok{\textbackslash{}n}\CommentTok{aI2}\CharTok{\textbackslash{}n}\CommentTok{aI3}\CharTok{\textbackslash{}n}\CommentTok{a.\textquotesingle{}}
\end{Highlighting}
\end{Shaded}

Le format n\textquotesingle est pas évident pour les lecteurs humains ;
il est conçu pour être facilement interprété par \texttt{pickle}.
\texttt{pickle.loads} (« load string ») reconstitue
l\textquotesingle objet{[}cite: 1{]} :

\begin{Shaded}
\begin{Highlighting}[]
\OperatorTok{\textgreater{}\textgreater{}\textgreater{}}\NormalTok{ t1 }\OperatorTok{=}\NormalTok{ [}\DecValTok{1}\NormalTok{, }\DecValTok{2}\NormalTok{, }\DecValTok{3}\NormalTok{]}
\OperatorTok{\textgreater{}\textgreater{}\textgreater{}}\NormalTok{ s }\OperatorTok{=}\NormalTok{ pickle.dumps(t1)}
\OperatorTok{\textgreater{}\textgreater{}\textgreater{}}\NormalTok{ t2 }\OperatorTok{=}\NormalTok{ pickle.loads(s)}
\OperatorTok{\textgreater{}\textgreater{}\textgreater{}} \BuiltInTok{print}\NormalTok{ t2}
\NormalTok{[}\DecValTok{1}\NormalTok{, }\DecValTok{2}\NormalTok{, }\DecValTok{3}\NormalTok{]}
\end{Highlighting}
\end{Shaded}

Bien que le nouvel objet ait la même valeur que
l\textquotesingle ancien, ce n\textquotesingle est (en général) pas le
même objet{[}cite: 1{]} :

\begin{Shaded}
\begin{Highlighting}[]
\OperatorTok{\textgreater{}\textgreater{}\textgreater{}}\NormalTok{ t1 }\OperatorTok{==}\NormalTok{ t2}
\VariableTok{True}
\OperatorTok{\textgreater{}\textgreater{}\textgreater{}}\NormalTok{ t1 }\KeywordTok{is}\NormalTok{ t2}
\VariableTok{False}
\end{Highlighting}
\end{Shaded}

En d\textquotesingle autres termes, sérialiser puis désérialiser
(\emph{pickling and unpickling}) a le même effet que copier
l\textquotesingle objet{[}cite: 1{]}.

Vous pouvez utiliser \texttt{pickle} pour stocker des types autres que
des chaînes dans une base de données. En fait, cette combinaison est si
courante qu\textquotesingle elle a été encapsulée dans un module appelé
\texttt{shelve}{[}cite: 1{]}.

\textbf{Exercice}

Si vous téléchargez ma solution à l\textquotesingle Exercice 4 depuis
\url{http://thinkpython.com/code/anagram_sets.py}, vous verrez
qu\textquotesingle elle crée un dictionnaire qui associe une chaîne
triée de lettres à la liste des mots qui peuvent être épelés avec ces
lettres. Par exemple, \texttt{\textquotesingle{}opst\textquotesingle{}}
est associé à la liste
\texttt{{[}\textquotesingle{}opts\textquotesingle{},\ \textquotesingle{}post\textquotesingle{},\ \textquotesingle{}pots\textquotesingle{},\ \textquotesingle{}spot\textquotesingle{},\ \textquotesingle{}stop\textquotesingle{},\ \textquotesingle{}tops\textquotesingle{}{]}}{[}cite:
1{]}. Écrivez un module qui importe \texttt{anagram\_sets} et fournit
deux nouvelles fonctions : \texttt{store\_anagrams} devrait stocker le
dictionnaire d\textquotesingle anagrammes dans un « shelf » ;
\texttt{read\_anagrams} devrait rechercher un mot et renvoyer la liste
de ses anagrammes{[}cite: 1{]}. Solution :
\url{http://thinkpython.com/code/anagram_db.py\%5Bcite}: 1{]}.

\hypertarget{pipes}{%
\subsection{Pipes}\label{pipes}}

La plupart des systèmes d\textquotesingle exploitation fournissent une
interface en ligne de commande, également connue sous le nom de
\emph{shell}. Les shells fournissent généralement des commandes pour
naviguer dans le système de fichiers et lancer des applications. Par
exemple, sous Unix, vous pouvez changer de répertoire avec \texttt{cd},
afficher le contenu d\textquotesingle un répertoire avec \texttt{ls}, et
lancer un navigateur web en tapant (par exemple)
\texttt{firefox}{[}cite: 1{]}.

N\textquotesingle importe quel programme que vous pouvez lancer depuis
le shell peut également être lancé depuis Python en utilisant un
\textbf{pipe}. Un pipe est un objet qui représente un programme en cours
d\textquotesingle exécution{[}cite: 1{]}.

Par exemple, la commande Unix \texttt{ls\ -l} affiche normalement le
contenu du répertoire courant (en format long). Vous pouvez lancer
\texttt{ls} avec \texttt{os.popen}{[}cite: 1{]} :

\begin{Shaded}
\begin{Highlighting}[]
\OperatorTok{\textgreater{}\textgreater{}\textgreater{}}\NormalTok{ cmd }\OperatorTok{=} \StringTok{\textquotesingle{}ls {-}l\textquotesingle{}}
\OperatorTok{\textgreater{}\textgreater{}\textgreater{}}\NormalTok{ fp }\OperatorTok{=}\NormalTok{ os.popen(cmd)}
\end{Highlighting}
\end{Shaded}

L\textquotesingle argument est une chaîne qui contient une commande
shell. La valeur de retour est un objet qui se comporte exactement comme
un fichier ouvert. Vous pouvez lire la sortie du processus \texttt{ls}
une ligne à la fois avec \texttt{readline} ou récupérer
l\textquotesingle ensemble en une seule fois avec \texttt{read}{[}cite:
1{]} :

\begin{Shaded}
\begin{Highlighting}[]
\OperatorTok{\textgreater{}\textgreater{}\textgreater{}}\NormalTok{ res }\OperatorTok{=}\NormalTok{ fp.read()}
\end{Highlighting}
\end{Shaded}

Lorsque vous avez terminé, vous fermez le pipe comme un fichier{[}cite:
1{]} :

\begin{Shaded}
\begin{Highlighting}[]
\OperatorTok{\textgreater{}\textgreater{}\textgreater{}}\NormalTok{ stat }\OperatorTok{=}\NormalTok{ fp.close()}
\OperatorTok{\textgreater{}\textgreater{}\textgreater{}} \BuiltInTok{print}\NormalTok{ stat}
\VariableTok{None}
\end{Highlighting}
\end{Shaded}

La valeur de retour est le statut final du processus \texttt{ls} ;
\texttt{None} signifie qu\textquotesingle il s\textquotesingle est
terminé normalement (sans erreurs){[}cite: 1{]}.

Par exemple, la plupart des systèmes Unix fournissent une commande
appelée \texttt{md5sum} qui lit le contenu d\textquotesingle un fichier
et calcule une « somme de contrôle » (\emph{checksum}). Vous pouvez lire
à propos de MD5 sur \url{http://en.wikipedia.org/wiki/Md5}. Cette
commande fournit un moyen efficace de vérifier si deux fichiers ont le
même contenu. La probabilité que des contenus différents produisent la
même somme de contrôle est très faible (c\textquotesingle est-à-dire peu
probable de se produire avant la fin de
l\textquotesingle univers){[}cite: 1{]}.

Vous pouvez utiliser un pipe pour exécuter \texttt{md5sum} depuis Python
et obtenir le résultat{[}cite: 1{]} :

\begin{Shaded}
\begin{Highlighting}[]
\OperatorTok{\textgreater{}\textgreater{}\textgreater{}}\NormalTok{ filename }\OperatorTok{=} \StringTok{\textquotesingle{}book.tex\textquotesingle{}}
\OperatorTok{\textgreater{}\textgreater{}\textgreater{}}\NormalTok{ cmd }\OperatorTok{=} \StringTok{\textquotesingle{}md5sum \textquotesingle{}} \OperatorTok{+}\NormalTok{ filename}
\OperatorTok{\textgreater{}\textgreater{}\textgreater{}}\NormalTok{ fp }\OperatorTok{=}\NormalTok{ os.popen(cmd)}
\OperatorTok{\textgreater{}\textgreater{}\textgreater{}}\NormalTok{ res }\OperatorTok{=}\NormalTok{ fp.read()}
\OperatorTok{\textgreater{}\textgreater{}\textgreater{}}\NormalTok{ stat }\OperatorTok{=}\NormalTok{ fp.close()}
\OperatorTok{\textgreater{}\textgreater{}\textgreater{}} \BuiltInTok{print}\NormalTok{ res}
\FloatTok{1e0033}\ErrorTok{f0ed0656636de0d75144ba32e0}\NormalTok{  book.tex}
\OperatorTok{\textgreater{}\textgreater{}\textgreater{}} \BuiltInTok{print}\NormalTok{ stat}
\VariableTok{None}
\end{Highlighting}
\end{Shaded}

\textbf{Exercice}

Dans une large collection de fichiers MP3, il peut y avoir plus
d\textquotesingle une copie de la même chanson, stockée dans des
répertoires différents ou avec des noms de fichiers différents. Le but
de cet exercice est de rechercher les doublons{[}cite: 1{]}. Écrivez un
programme qui recherche dans un répertoire et tous ses sous-répertoires,
de manière récursive, et renvoie une liste de chemins complets pour tous
les fichiers avec un suffixe donné (comme \texttt{.mp3}){[}cite: 1{]}.
\emph{Indice :} \texttt{os.path} fournit plusieurs fonctions utiles pour
manipuler les noms de fichiers et de chemins{[}cite: 1{]}. Pour
reconnaître les doublons, vous pouvez utiliser \texttt{md5sum} pour
calculer une somme de contrôle pour chaque fichier. Si deux fichiers ont
la même somme de contrôle, ils ont probablement le même contenu. Pour
vérifier à deux reprises, vous pouvez utiliser la commande Unix
\texttt{diff}{[}cite: 1{]}. Solution :
\url{http://thinkpython.com/code/find_duplicates.py\%5Bcite}: 1{]}.

\hypertarget{uxe9criture-de-modules}{%
\subsection{Écriture de modules}\label{uxe9criture-de-modules}}

N\textquotesingle importe quel fichier qui contient du code Python peut
être importé en tant que module. Par exemple, supposons que vous ayez un
fichier nommé \texttt{wc.py} avec le code suivant{[}cite: 1{]} :

\begin{Shaded}
\begin{Highlighting}[]
\KeywordTok{def}\NormalTok{ linecount(filename):}
\NormalTok{    count }\OperatorTok{=} \DecValTok{0}
    \ControlFlowTok{for}\NormalTok{ line }\KeywordTok{in} \BuiltInTok{open}\NormalTok{(filename):}
\NormalTok{        count }\OperatorTok{+=} \DecValTok{1}
    \ControlFlowTok{return}\NormalTok{ count}

\BuiltInTok{print}\NormalTok{ linecount(}\StringTok{\textquotesingle{}wc.py\textquotesingle{}}\NormalTok{)}
\end{Highlighting}
\end{Shaded}

Si vous exécutez ce programme, il se lit lui-même et affiche le nombre
de lignes dans le fichier, qui est 7{[}cite: 1{]}.

Vous pouvez également l\textquotesingle importer de cette façon{[}cite:
1{]} :

\begin{Shaded}
\begin{Highlighting}[]
\OperatorTok{\textgreater{}\textgreater{}\textgreater{}} \ImportTok{import}\NormalTok{ wc}
\DecValTok{7}
\end{Highlighting}
\end{Shaded}

Vous disposez désormais d\textquotesingle un objet module
\texttt{wc}{[}cite: 1{]} :

\begin{Shaded}
\begin{Highlighting}[]
\OperatorTok{\textgreater{}\textgreater{}\textgreater{}} \BuiltInTok{print}\NormalTok{ wc}
\OperatorTok{\textless{}}\NormalTok{module }\StringTok{\textquotesingle{}wc\textquotesingle{}} \ImportTok{from} \StringTok{\textquotesingle{}wc.py\textquotesingle{}}\OperatorTok{\textgreater{}}
\end{Highlighting}
\end{Shaded}

Qui fournit une fonction appelée \texttt{linecount}{[}cite: 1{]} :

\begin{Shaded}
\begin{Highlighting}[]
\OperatorTok{\textgreater{}\textgreater{}\textgreater{}}\NormalTok{ wc.linecount(}\StringTok{\textquotesingle{}wc.py\textquotesingle{}}\NormalTok{)}
\DecValTok{7}
\end{Highlighting}
\end{Shaded}

C\textquotesingle est donc ainsi que l\textquotesingle on écrit des
modules en Python{[}cite: 1{]}.

Le seul problème avec cet exemple est que lorsque vous importez le
module, il exécute le code de test en bas. Normalement, lorsque vous
importez un module, il définit de nouvelles fonctions mais il ne les
exécute pas{[}cite: 1{]}.

Les programmes qui seront importés comme modules utilisent souvent
l\textquotesingle idiome suivant{[}cite: 1{]} :

\begin{Shaded}
\begin{Highlighting}[]
\ControlFlowTok{if} \VariableTok{\_\_name\_\_} \OperatorTok{==} \StringTok{\textquotesingle{}\_\_main\_\_\textquotesingle{}}\NormalTok{:}
    \BuiltInTok{print}\NormalTok{ linecount(}\StringTok{\textquotesingle{}wc.py\textquotesingle{}}\NormalTok{)}
\end{Highlighting}
\end{Shaded}

\texttt{\_\_name\_\_} est une variable intégrée qui est définie au
démarrage du programme. Si le programme s\textquotesingle exécute en
tant que script, \texttt{\_\_name\_\_} a la valeur \texttt{\_\_main\_\_}
; dans ce cas, le code de test est exécuté. Sinon, si le module est en
cours d\textquotesingle importation, le code de test est ignoré{[}cite:
1{]}.

\textbf{Exercice}

Tapez cet exemple dans un fichier nommé \texttt{wc.py} et exécutez-le en
tant que script. Ensuite, lancez l\textquotesingle interpréteur Python
et faites \texttt{import\ wc}. Quelle est la valeur de
\texttt{\_\_name\_\_} lorsque le module est en cours
d\textquotesingle importation ?{[}cite: 1{]} \emph{Avertissement :} Si
vous importez un module qui a déjà été importé, Python ne fait rien. Il
ne relit pas le fichier, même s\textquotesingle il a changé. Si vous
souhaitez recharger un module, vous pouvez utiliser la fonction intégrée
\texttt{reload}, mais cela peut être délicat, donc la chose la plus sûre
à faire est de redémarrer l\textquotesingle interpréteur puis
d\textquotesingle importer à nouveau le module{[}cite: 1{]}.

\hypertarget{duxe9bogage}{%
\subsection{Débogage}\label{duxe9bogage}}

Lorsque vous lisez et écrivez des fichiers, vous pourriez rencontrer des
problèmes avec les espaces blancs. Ces erreurs peuvent être difficiles à
déboguer parce que les espaces, les tabulations et les retours à la
ligne sont normalement invisibles{[}cite: 1{]} :

\begin{Shaded}
\begin{Highlighting}[]
\OperatorTok{\textgreater{}\textgreater{}\textgreater{}}\NormalTok{ s }\OperatorTok{=} \StringTok{\textquotesingle{}1 2}\CharTok{\textbackslash{}t}\StringTok{ 3}\CharTok{\textbackslash{}n}\StringTok{ 4\textquotesingle{}}
\OperatorTok{\textgreater{}\textgreater{}\textgreater{}} \BuiltInTok{print}\NormalTok{ s}
\DecValTok{1} \DecValTok{2}  \DecValTok{3}
 \DecValTok{4}
\end{Highlighting}
\end{Shaded}

La fonction intégrée \texttt{repr} peut aider. Elle prend
n\textquotesingle importe quel objet comme argument et renvoie une
représentation sous forme de chaîne de l\textquotesingle objet. Pour les
chaînes, elle représente les caractères d\textquotesingle espacement
avec des séquences d\textquotesingle antislash{[}cite: 1{]} :

\begin{Shaded}
\begin{Highlighting}[]
\OperatorTok{\textgreater{}\textgreater{}\textgreater{}} \BuiltInTok{print} \BuiltInTok{repr}\NormalTok{(s)}
\CommentTok{\textquotesingle{}1 2}\CharTok{\textbackslash{}t}\CommentTok{ 3}\CharTok{\textbackslash{}n}\CommentTok{ 4\textquotesingle{}}
\end{Highlighting}
\end{Shaded}

Cela peut être utile pour le débogage{[}cite: 1{]}.

Un autre problème que vous pourriez rencontrer est que différents
systèmes utilisent des caractères différents pour indiquer la fin
d\textquotesingle une ligne. Certains systèmes utilisent un saut de
ligne, représenté par \texttt{\textbackslash{}n}.
D\textquotesingle autres utilisent un caractère de retour, représenté
par \texttt{\textbackslash{}r}. Certains utilisent les deux{[}cite:
1{]}.

Si vous déplacez des fichiers entre différents systèmes, ces
incohérences pourraient causer des problèmes. Pour la plupart des
systèmes, il existe des applications pour convertir d\textquotesingle un
format à un autre. Vous pouvez les trouver (et en lire plus sur ce
problème) à l\textquotesingle adresse
\url{http://en.wikipedia.org/wiki/Newline}. Ou, bien sûr, vous pourriez
en écrire un vous-même{[}cite: 1{]}.

\hypertarget{glossaire}{%
\subsection{Glossaire}\label{glossaire}}

persistant persistent Relatif à un programme qui
s\textquotesingle exécute indéfiniment et conserve au moins une partie
de ses données dans un stockage permanent{[}cite: 1{]}.

opérateur de formatage format operator Un opérateur, \texttt{\%}, qui
prend une chaîne de format et un tuple et génère une chaîne incluant les
éléments du tuple formatés tel que spécifié par la chaîne de
format{[}cite: 1{]}.

chaîne de format format string Une chaîne, utilisée avec
l\textquotesingle opérateur de formatage, qui contient des séquences de
format{[}cite: 1{]}.

séquence de format format sequence Une séquence de caractères dans une
chaîne de format, comme \texttt{\%d}, qui spécifie comment une valeur
doit être formatée{[}cite: 1{]}.

fichier texte text file Une séquence de caractères stockée dans un
stockage permanent comme un disque dur{[}cite: 1{]}.

répertoire directory Une collection nommée de fichiers, également
appelée dossier{[}cite: 1{]}.

chemin path Une chaîne qui identifie un fichier{[}cite: 1{]}.

chemin relatif relative path Un chemin qui part du répertoire
courant{[}cite: 1{]}.

chemin absolu absolute path Un chemin qui part du répertoire le plus
haut du système de fichiers{[}cite: 1{]}.

capturer catch Empêcher une exception de terminer un programme en
utilisant les instructions \texttt{try} et \texttt{except}{[}cite: 1{]}.

base de données database Un fichier dont le contenu est organisé comme
un dictionnaire avec des clés qui correspondent à des valeurs{[}cite:
1{]}.

\hypertarget{exercices}{%
\subsection{Exercices}\label{exercices}}

\textbf{Exercice}

Le module \texttt{urllib} fournit des méthodes pour manipuler les URL et
télécharger des informations depuis le web. Le code suivant télécharge
et affiche un message secret provenant de thinkpython.com{[}cite: 1{]} :

\begin{Shaded}
\begin{Highlighting}[]
\ImportTok{import}\NormalTok{ urllib}

\NormalTok{conn }\OperatorTok{=}\NormalTok{ urllib.urlopen(}\StringTok{\textquotesingle{}http://thinkpython.com/secret.html\textquotesingle{}}\NormalTok{)}
\ControlFlowTok{for}\NormalTok{ line }\KeywordTok{in}\NormalTok{ conn:}
    \BuiltInTok{print}\NormalTok{ line.strip()}
\end{Highlighting}
\end{Shaded}

Exécutez ce code et suivez les instructions que vous y voyez{[}cite:
1{]}. Solution : \url{http://thinkpython.com/code/zip_code.py\%5Bcite}:
1{]}.
